\section{Conclusiones}
% ==========================================
La biblioteca \texttt{ftir\_library} v0.0.4 consolida un entorno unificado y maduro para la caracterización, control de calidad y diagnóstico asistido por computadora utilizando espectroscopía vibracional (FTIR y Raman SERS). 

El desacoplamiento de la arquitectura a través del módulo \path{stats.py} permite reutilizar de manera eficiente algoritmos avanzados como la proyección de componentes principales (PCA) mediante SVD y clasificadores regularizados multiclase (Kernel Ridge Classifier), habilitando flujos de trabajo científicos flexibles. El soporte nativo para el estándar industrial JCAMP-DX expande significativamente la utilidad de la biblioteca al permitir la ingesta directa de bases de datos de referencia clásicas como la del NIST. 

Finalmente, la integración de una suite completa de pruebas automatizadas mediante \texttt{pytest} garantiza la estabilidad numérica y el correcto funcionamiento del pipeline frente a futuras extensiones, estableciendo una base sólida para el desarrollo de herramientas de diagnóstico biomédico y caracterización de nanomateriales en entornos de producción.
\newpage